% libre-compendium
%
% Copyright (C) 2026 Richard Qin
%
% This document is dual-licensed under either:
% - GNU Free Documentation License, Version 1.3 or any later version
%   published by the Free Software Foundation, with no Invariant Sections,
%   no Front-Cover Texts, and no Back-Cover Texts (see COPYING.DOC); or
% - Creative Commons Attribution-ShareAlike 4.0 International License
%   (see LICENSE-CC-BY-SA-4.0).
\ProvidesFile{settings.tex}[2026/11/15 v0.1.0 离散数学设置]
\lcset{
  course           = {离散数学},
  copyright-holder = {Richard Qin},
  copyright-years  = {2026},
}
\lchistory{2026-06-10}{v0.1.0}{Richard Qin}{完成集合、关系、逻辑与代数结构三份讲义初稿}
\lchistory{2026-11-15}{v0.2.0}{Richard Qin}{合并集合、关系、逻辑与代数结构三份讲义并迁移至 common}
\allowdisplaybreaks
\colorlet{CADef}{lc-def}
\colorlet{CAThm}{lc-thm}
\colorlet{CAMethod}{lc-method}
\colorlet{CAPitfall}{lc-pitfall}
\colorlet{CAExample}{lc-example}
\newcommand{\card}[1]{\lvert#1\rvert}
\newcommand{\pow}{\mathcal P}
\newcommand{\rel}{\mathrel{R}}
\newcommand{\id}{I}
\newcommand{\dom}{\operatorname{dom}}
\newcommand{\ran}{\operatorname{ran}}
\newcommand{\lub}{\operatorname{lub}}
\newcommand{\glb}{\operatorname{glb}}
\newcommand{\ord}{\operatorname{ord}}
\newcommand{\Ker}{\operatorname{Ker}}
\newcommand{\img}{\operatorname{Im}}
\newcommand{\lcm}{\operatorname{lcm}}
\newcommand{\T}{\mathbf{1}}
\newcommand{\F}{\mathbf{0}}
\newcommand{\wff}{\mathrm{wff}}
\newcommand{\vDashs}{\vDash}
\newcommand{\equivv}{\Leftrightarrow}
